\begin{python0}
from solutions import *; clear()
\end{python0}
\sectionthree{Default Values}

Try this:

\begin{consolethree}[escapeinside=||]
#include <iostream>

void printHeads(int numHeads = 1)
{
    std::cout << "number of heads: " << numHeads
              << std::endl;
}

int main()
{
    printHeads(); // NO ARGUMENT PASSED IN!!!
    printHeads(42);
    return 0;
}
\end{consolethree}

(Forget about prototypes for the time being. They will come in a bit ...)

Parameter \texttt{numHeads} of function \texttt{printHeads()} is called a parameter with a \EMPHASIZE{default value}.

One obvious use of default values is that you can make it easier for other programmers to use your function: You choose the most common value for a parameter and make it the default value for that parameter. Any time someone (including yourself) uses the function in the ``most commonly used'' manner, he/she need not send in an argument.

Here's another example:

\begin{consolethree}[escapeinside=||]
#include <iostream>

void printHello(int i = 0)
{
    std::cout << "hello ";
    switch (i)
    {
        case 0: std::cout << "world\n"; break;
        case 1: std::cout << "galaxy\n"; break;
        case 2: std::cout << "universe\n"; break;
        case 3: std::cout << "underverse\n"; break;
    }
}

int main()
{
    printHello();
    printHello(1);
    return 0;
}
\end{consolethree}

Right now there are only \EMPHASIZE{three rules} to remember: Think of the parameters of a function as parameters without default values and parameters with default values. \EMPHASIZE{All the parameters with default values must be together} and appear \EMPHASIZE{after those without default values}.

In other words this is OK:

\begin{consolethree}[escapeinside=||]
bool spam(int a, double b, |\textbf{\underline{char c = '\$', int d = 42}}|)
{
    ...
}
\end{consolethree}

But this is WRONG:

\begin{consolethree}[escapeinside=||]
bool spam(int a, |\textbf{\underline{double b = 3.14}}|, char c, |\textbf{\underline{int d = 42}}|)|\xsidebox[0cm]{0cm}{}{txt2}{\texttt{c} does not have a default value but appears after \texttt{b} which has a default value.}|
{
    ...
}
\end{consolethree}

Make sure you try this or you won't remember.

So remember to group up the parameters into those with default values
and those without. The ones without (if any) must come first.

Another rule: \textbf{default values must be constants} (example:
constant literals such as 5 or constant expressions.) In other words
this is OK:

\begin{consolethree}[escapeinside=||]
const int X = 42;
bool spam(int a, double b, char c = '$', int d = X+1)
{
    ...
}
\end{consolethree}

But this is WRONG:

\begin{consolethree}[escapeinside=||]
int X = 42;
bool spam(int a, double b, char c = '$', int d = X)
{
    ...
}
\end{consolethree}

\begin{ex}
What is the output (or circle the errors)? (Do this by hand)

\begin{consolethree}[escapeinside=||]
int f(int a, int b = 4, int c = 5)
{
    return a + b + c;
}

int main()
{
    std::cout << f(1, 2, 3) << ' ' << f(1, 2) << ' '
              << f(1) << std::endl;
    return 0;
}
\end{consolethree}

Now verify using your compiler.
\end{ex}

Last one: \EMPHASIZE{You cannot have default values for array parameters.} For instance the following will not compile:

\begin{consolethree}[escapeinside=||]
#include <iostream>

void print(int x[] = {2, 3})
{
    std::cout << x[0] << ' ' << x[1];
}

int main()
{
    print();
    return 0;
}
\end{consolethree}

(See last section for exception.)

\begin{ex}
What is the output (or circle the errors)? (Do this by hand)

\begin{consolethree}[escapeinside=||]
int f(int a = 3, int b, int c = 5)
{
    return a + b + c;
}

int main()
{
    std::cout << f(1, 2, 3) << std::endl;
    return 0;
}
\end{consolethree}

Now verify using your compiler.
\end{ex}

\begin{ex}
Write a function \texttt{printDateTime()} that prints the date and time data passed in:

\begin{consolethree}[escapeinside=||]
void printDateTime(int yyyy, int mm, int dd,
                   int hh = 0, int mm = 0, int ss = 0)
{
    ... CODE ...
}

int main()
{
    printDateTime(2007, 1, 15);
    printDateTime(2007, 1, 15, 14, 15, 59);
    return 0;
}
\end{consolethree}

Expected output:

\begin{console}
2007-01-15 00:00:00
2007-01-15 14:15:59
\end{console}

After you are done, make January 1, 1970 the default date.

[Hint: Use \texttt{std::fill()} function.]
\end{ex}


